\chapter*{Programa de la Materia}

\section*{Unidad 1: Introducción}

\paragraph*{Sistemas de cómputos} Elementos básicos. Ejecución de intrucciones. Interrupciones. Jerarquía de memorias. Técnicas de comunicación de E/S.

\paragraph*{Sistemas operativos} Funciones y objetivos de los sistemas operativos. Estructura de los sistemas operativos. Características de los sistemas modernos. Consideraciones prácticas para la selección del sistema según el entorno de aplicación.

\section*{Unidad 2: Procesos}

\paragraph*{Descripción y control de procesos e hilos} Estados de un proceso. Estructuras de datos y control de procesos. Concepto de procesos e hilos. Estados de un hilo. Implementación de hilos.

\paragraph*{Planificación} Niveles de planificación. Objetivos y criterios. Planificación de monoprocesadores: prioridades, estrategias, algoritmos de planificación. Planificación de multiprocesadores. Planificación en sistemas de tiempo real.

\section*{Unidad 3: Concurrencia}

\paragraph*{Exclusión mutua y sincronización} Principios generales de concurrencia. Soluciones por software. Soluciones por hardware. Semáforos. El problema del productor consumidor. Monitores. Paso de mensajes.

\paragraph*{Bloqueo mutuo} Principios del bloqueo mutuo. Prevención del bloqueo mutuo. Predicción del bloqueo mutuo. Detección del bloqueo mutuo. Estrategias integradas.

\section*{Unidad 4: Memoria}

\paragraph*{Gestión de memoria} Requisitos de la gestión de la memoria. Particiones fijas y dinámicas. Estrategias de ubicación. Paginación. Segmentación.

\paragraph*{Memoria virtual} Principio de cercanía. Paginado y segmentado con memoria virtual. Estrategias de gestión de memoria virtual.

\section*{Unidad 5: Entrada salida y archivos}

\paragraph*{Gestión de E/S y planificación de discos} Organización de las funciones de E/S. Aspectos de diseño. Buffering. Planificación de discos. RAID.

\paragraph*{Sistemas de archivos} Características generales. Organización y acceso a archivos. Directorios. Agrupación de registros. Gestión del almacenamiento secundario.

\section*{Unidad 6: Virtualización y contenerización}

\paragraph*{Introducción a la virtualización} Conceptos. Emulación. Virtualización Total. Paravirtualización. Contenedores. Portabilidad. Usos comunes.

\paragraph*{Virtualización} Virtualización de hardware. Virtualización de red. Virtualización de almacenamiento. Soporte de hardware para virtualizar. Implementaciones: Proxmox - KVM. VirtualBox.

\paragraph*{Contenedores} Conceptos. Virtualización a nivel de SO. Arquitectura. Portabilidad. Implementaciones: Proxmox - LXC. Docker.

\section*{Unidad 7: Conceptos básicos de administración y operación}

\paragraph*{Linux} Características generales. Organización de directorios. Shells. Comandos. Editores. Scripts. Manejo de usuarios, permisos y archivos. Procesos. Servicios del sistema. Tareas de administración.

\paragraph*{Windows} Características generales. Redes basadas en Microsoft Windows. Conceptos de dominios y directorio activo. Manejo de usuarios, grupos, UO. Servicios del sistema. Sistema de archivos NTFS. Permisos, derechos, herencia. Tareas de administración.

